\documentclass[11pt]{article}

\usepackage[margin=1in]{geometry}
\usepackage{amsmath,amssymb,bm}
\usepackage{array}
\usepackage{booktabs}
\usepackage{xcolor}
\usepackage{tcolorbox}
\usepackage{enumitem}

\title{From the General Elastic Stiffness Tensor to VTI and TTI Media}
\author{}
\date{}

\begin{document}

\maketitle

\section{The general stress--strain relation}

For a linear elastic material, stress and strain are related by Hooke's law,

\begin{equation}
    \sigma_{ij} = C_{ijkl}\,\epsilon_{kl},
\end{equation}

where

\begin{itemize}
    \item $\sigma_{ij}$ is the stress tensor,
    \item $\epsilon_{kl}$ is the strain tensor,
    \item $C_{ijkl}$ is the fourth-order stiffness tensor.
\end{itemize}

In three dimensions, each index can take the values $1$, $2$, or $3$. Therefore, at first sight, the stiffness tensor contains

\begin{equation}
    3^4 = 81
\end{equation}

coefficients.

However, most of these coefficients are not independent.

\section{From 81 coefficients to 36: minor symmetries}

For ordinary elastic materials, both the stress tensor and the strain tensor are symmetric:

\begin{equation}
    \sigma_{ij} = \sigma_{ji},
\end{equation}

and

\begin{equation}
    \epsilon_{ij} = \epsilon_{ji}.
\end{equation}

These symmetries imply the so-called \emph{minor symmetries} of the stiffness tensor:

\begin{equation}
    C_{ijkl} = C_{jikl},
\end{equation}

and

\begin{equation}
    C_{ijkl} = C_{ijlk}.
\end{equation}

Thus, instead of the pair $ij$ having nine possible combinations,

\[
11,\;12,\;13,\;21,\;22,\;23,\;31,\;32,\;33,
\]

only six are distinct:

\[
11,\quad 22,\quad 33,\quad 23,\quad 13,\quad 12.
\]

The same is true for the pair $kl$.

Therefore, the fourth-order tensor can be represented as a $6\times 6$ matrix:

\begin{equation}
    81 \longrightarrow 36.
\end{equation}

\section{Voigt notation}

The six independent index pairs are commonly replaced by a single index:

\begin{equation}
\begin{aligned}
    11 &\rightarrow 1,\\
    22 &\rightarrow 2,\\
    33 &\rightarrow 3,\\
    23 &\rightarrow 4,\\
    13 &\rightarrow 5,\\
    12 &\rightarrow 6.
\end{aligned}
\end{equation}

Then the stiffness tensor $C_{ijkl}$ becomes the matrix $C_{IJ}$, and the stress--strain relation may be written schematically as

\begin{equation}
\begin{pmatrix}
\sigma_{11}\\
\sigma_{22}\\
\sigma_{33}\\
\sigma_{23}\\
\sigma_{13}\\
\sigma_{12}
\end{pmatrix}
=
\begin{pmatrix}
C_{11}&C_{12}&C_{13}&C_{14}&C_{15}&C_{16}\\
C_{21}&C_{22}&C_{23}&C_{24}&C_{25}&C_{26}\\
C_{31}&C_{32}&C_{33}&C_{34}&C_{35}&C_{36}\\
C_{41}&C_{42}&C_{43}&C_{44}&C_{45}&C_{46}\\
C_{51}&C_{52}&C_{53}&C_{54}&C_{55}&C_{56}\\
C_{61}&C_{62}&C_{63}&C_{64}&C_{65}&C_{66}
\end{pmatrix}
\begin{pmatrix}
\epsilon_{11}\\
\epsilon_{22}\\
\epsilon_{33}\\
\epsilon_{23}\\
\epsilon_{13}\\
\epsilon_{12}
\end{pmatrix}.
\end{equation}

The exact factors of two associated with the shear strains depend on the convention used for Voigt notation, but the important point here is that the fourth-order tensor can be represented by a $6\times6$ matrix.

\section{From 36 coefficients to 21: major symmetry}

For a conventional elastic material, the elastic energy must be independent of the order in which stresses and strains are applied. This leads to the \emph{major symmetry}

\begin{equation}
    C_{ijkl} = C_{klij}.
\end{equation}

In Voigt notation, this becomes simply

\begin{equation}
    C_{IJ} = C_{JI}.
\end{equation}

Therefore, the $6\times6$ stiffness matrix is symmetric.

A symmetric $6\times6$ matrix contains

\begin{equation}
    \frac{6(6+1)}{2} = 21
\end{equation}

independent coefficients.

Thus,

\begin{equation}
    \boxed{81 \longrightarrow 36 \longrightarrow 21}.
\end{equation}

\begin{tcolorbox}[title=Interpretation]
The reductions from 81 to 36 and then to 21 do not yet assume any special geological symmetry of the material.

They follow from the symmetry of stress and strain and from the existence of an elastic strain-energy function.
\end{tcolorbox}

The most general linear elastic material therefore requires 21 independent stiffness coefficients. Such a material belongs to the \emph{triclinic} symmetry class.

\section{Material symmetry}

The next reductions come from assuming that the material itself has symmetry.

The basic idea is:

\begin{quote}
If the material behaves identically after certain rotations or reflections, some stiffness coefficients must be equal and others must vanish.
\end{quote}

The greater the symmetry of the material, the fewer independent elastic constants are required.

At one extreme, a completely isotropic material looks mechanically identical in every direction and requires only two independent elastic constants, for example the Lam\'e parameters

\begin{equation}
    \lambda,\qquad \mu.
\end{equation}

Thus,

\begin{equation}
    21 \longrightarrow 2
\end{equation}

for an isotropic elastic medium.

\section{Transverse isotropy}

A transversely isotropic material has one special direction, called the \emph{symmetry axis}.

All directions perpendicular to this axis are mechanically equivalent.

A good geological example is a finely layered sedimentary medium. The medium can behave differently parallel and perpendicular to the layering, while directions within the plane of the layers are equivalent.

For a transversely isotropic medium, only five independent stiffness coefficients are required.

\section{VTI: vertical transverse isotropy}

For a VTI medium, the symmetry axis is vertical. If the vertical direction is $x_3=z$, then the horizontal directions $x_1$ and $x_2$ are equivalent.

The stiffness matrix becomes

\begin{equation}
\mathbf C_{\mathrm{VTI}}
=
\begin{pmatrix}
C_{11} & C_{12} & C_{13} & 0 & 0 & 0\\
C_{12} & C_{11} & C_{13} & 0 & 0 & 0\\
C_{13} & C_{13} & C_{33} & 0 & 0 & 0\\
0 & 0 & 0 & C_{44} & 0 & 0\\
0 & 0 & 0 & 0 & C_{44} & 0\\
0 & 0 & 0 & 0 & 0 & C_{66}
\end{pmatrix}.
\end{equation}

The coefficient $C_{12}$ is not independent because

\begin{equation}
    C_{12} = C_{11} - 2C_{66}.
\end{equation}

Therefore, a convenient set of five independent constants is

\begin{equation}
    \boxed{
    C_{11},\quad
    C_{33},\quad
    C_{13},\quad
    C_{44},\quad
    C_{66}.
    }
\end{equation}

Thus,

\begin{equation}
    \boxed{21 \longrightarrow 5}
\end{equation}

for transverse isotropy.

\subsection{Physical interpretation}

A useful approximate interpretation is:

\begin{itemize}
    \item $C_{33}$ mainly controls $P$-wave propagation along the vertical symmetry axis,
    \item $C_{11}$ mainly controls horizontal $P$-wave propagation,
    \item $C_{44}$ controls shear deformation involving the symmetry axis,
    \item $C_{66}$ controls shear deformation in the horizontal plane,
    \item $C_{13}$ couples vertical and horizontal deformation.
\end{itemize}

\section{Thomsen parameters}

In exploration geophysics, VTI media are often described using vertical velocities and Thomsen anisotropy parameters instead of the stiffness coefficients directly.

A common parameterization is

\begin{equation}
    V_{P0},\qquad
    V_{S0},\qquad
    \epsilon,\qquad
    \delta,\qquad
    \gamma,
\end{equation}

together with density $\rho$.

The vertical velocities are related to the stiffnesses by

\begin{equation}
    C_{33} = \rho V_{P0}^2,
\end{equation}

and

\begin{equation}
    C_{44} = \rho V_{S0}^2.
\end{equation}

The Thomsen parameter $\epsilon$ is

\begin{equation}
    \epsilon
    =
    \frac{C_{11}-C_{33}}{2C_{33}},
\end{equation}

so that

\begin{equation}
    C_{11}
    =
    C_{33}(1+2\epsilon).
\end{equation}

Similarly,

\begin{equation}
    \gamma
    =
    \frac{C_{66}-C_{44}}{2C_{44}},
\end{equation}

and therefore

\begin{equation}
    C_{66}
    =
    C_{44}(1+2\gamma).
\end{equation}

The parameter $\delta$ controls the near-vertical angular dependence of $P$-wave propagation and determines the remaining coupling stiffness $C_{13}$.

Thus, the Thomsen parameters provide a physically convenient description of the same five independent elastic constants of a VTI material.

\section{From VTI to TTI}

TTI stands for \emph{Tilted Transverse Isotropy}.

An important conceptual point is that TTI is not a new intrinsic material symmetry class. It is still a transversely isotropic material.

The difference is simply the orientation of the symmetry axis.

For VTI,

\begin{equation}
    \text{symmetry axis} = \text{vertical}.
\end{equation}

For TTI,

\begin{equation}
    \text{symmetry axis} = \text{tilted}.
\end{equation}

This is useful for describing dipping sedimentary layers.

In a local coordinate system aligned with the layers, the medium can first be described as VTI using the five independent stiffness coefficients

\begin{equation}
    C_{11},\quad C_{33},\quad C_{13},\quad C_{44},\quad C_{66}.
\end{equation}

The stiffness tensor is then rotated into the global computational coordinates.

If $R$ is the rotation matrix, the rotated fourth-order tensor is

\begin{equation}
    C^{\mathrm{TTI}}_{ijkl}
    =
    R_{im}R_{jn}R_{kp}R_{lq}
    C^{\mathrm{VTI}}_{mnpq}.
\end{equation}

For an introductory course, the main idea is more important than memorizing the full tensor rotation:

\begin{equation}
\boxed{
\text{Construct VTI in coordinates aligned with the layers}
\quad\longrightarrow\quad
\text{rotate the tensor into the global coordinates}.
}
\end{equation}

In a two-dimensional $x$--$z$ problem, the orientation is commonly specified by a single tilt angle $\theta$.

In three dimensions, two angles are generally required to specify the orientation of the symmetry axis, for example tilt and azimuth.

\section{Summary}

The progression from the completely general stiffness tensor to commonly used geophysical models can be summarized as

\begin{equation}
\boxed{
\begin{array}{rcl}
3^4=81
&\longrightarrow&
\text{general fourth-order tensor},
\\[2mm]
81\rightarrow36
&\longrightarrow&
\text{minor symmetries of stress and strain},
\\[2mm]
36\rightarrow21
&\longrightarrow&
\text{major symmetry from elastic energy},
\\[2mm]
21\rightarrow5
&\longrightarrow&
\text{transverse isotropy},
\\[2mm]
5
&\longrightarrow&
\text{VTI if the symmetry axis is vertical},
\\[2mm]
5+\text{orientation}
&\longrightarrow&
\text{TTI if the symmetry axis is tilted},
\\[2mm]
5\rightarrow2
&\longrightarrow&
\text{isotropy}.
\end{array}
}
\end{equation}

\begin{tcolorbox}[title=Key idea for students]
The 81 components arise because the stiffness is a fourth-order tensor in three dimensions.

Symmetries of stress, strain, and elastic energy reduce the number of independent coefficients to 21.

Material symmetry reduces the number further. A transversely isotropic material requires five independent elastic constants. If its symmetry axis is vertical, the medium is VTI. If the same symmetry axis is tilted, the medium is TTI.
\end{tcolorbox}

\end{document}
